\section{Excluding finite mass-shell profiles}
\label{sec:trace-profiles}
The centered trace has a short-distance behavior that can distinguish it from
finite polynomial expressions in free fields. The relevant distinction is the
negative power of the logarithm. We prove the exclusion directly for the
mass-shell profiles appearing in the statement. The proof proceeds through
finite power--logarithm approximations to Laplace integrals; it does not
assume a convergent asymptotic series.

\subsection{An algebra of finite asymptotic approximations}
Let $\mathcal P$ be the real vector space of finite sums of the functions
$s^a(\log s)^k$, where $a\in\mathbb Z$, $k\in\mathbb N$ and $s>0$.
It is an algebra containing the constants. Say that $f$ has an approximation
of order $b\in\mathbb Z$ if
\[
 f(s)=P(s)+O(s^b)\quad(s\downarrow0),\qquad P\in\mathcal P.
\]
Say that $f$ has power--logarithm approximations if it has an approximation
of every integer order. Different orders may use different finite sums.
An approximation of order $b$ is also one of any order $a\leq b$, since
$s^b\leq s^a$ for $0<s<1$.

\begin{lemma}\label{lem:profile-algebra}
Functions with power--logarithm approximations form a real algebra. Every
such function, and every element of $\mathcal P$, is $O(s^u)$ for some
integer $u$. Multiplication by $s^a$ shifts the order of an approximation
by $a$.
\end{lemma}
\begin{proof}
For $0<s<1$, $|\log s|\leq s^{-1}$. Thus each monomial is bounded by
$s^{a-k}$, and a finite sum is $O(s^u)$ for the smallest of the finitely
many exponents $a-k$. An order-zero approximation then gives the growth
claim for $f$, after replacing $u$ by $\min(u,0)$. The assertions about
addition, real scalar multiplication, and shifts are immediate from the
definition and closure of $\mathcal P$ under these operations.

For multiplication fix a desired order $b$. Choose $\ell$ with
$g=O(s^\ell)$ and an approximation $f=P+O(s^{b-\ell})$. Choose $u$ with
$P=O(s^u)$ and then $g=Q+O(s^{b-u})$. The exact identity
\[
 fg-PQ=(f-P)g+P(g-Q)
\]
shows that both terms are $O(s^b)$. Since $PQ\in\mathcal P$, this proves
the required approximation. In particular it proves closure under every
finite product, including the empty product $1$.
\end{proof}

\begin{lemma}[Integrating an approximation]\label{lem:profile-primitive}
Suppose $F$ is differentiable on $(0,\infty)$, $F'=g$, and $g$ has an
approximation of integer order $b\neq-1$. Then $F$ has an approximation
of order $b+1$.
\end{lemma}
\begin{proof}
Every monomial has an antiderivative in $\mathcal P$. For $a=-1$ it is
$(\log s)^{k+1}/(k+1)$. For $a\neq-1$, integration by parts gives the
recursion
\[
 I_{a,k}(s)=\frac{s^{a+1}(\log s)^k}{a+1}
       -\frac{k}{a+1}I_{a,k-1}(s),\qquad
 I_{a,0}(s)=\frac{s^{a+1}}{a+1}.
\]
These formulas are verified by differentiation, so no integrability at zero
is being assumed. If $g=P+O(s^b)$, subtract a primitive of $P$ from $F$.
It remains to treat a differentiable $R$ with
$|R'(s)|\leq Cs^b$ on $(0,\delta)$, where $C\geq0$.

Put $d=b+1\neq0$ and $A=C/d$. The derivatives of
$As^d+R(s)$ and $As^d-R(s)$ are nonnegative. The mean value theorem
therefore makes both functions increasing. Consequently, for
$0<x<y<\delta$,
\begin{equation}\label{eq:profile-comparison}
 |R(y)-R(x)|\leq A(y^d-x^d).
\end{equation}
The right-hand side is nonnegative also when $d<0$. In that case fix
$y\in(0,\delta)$; the inequality bounds $|R(x)|$ by a constant plus
$|A|x^d$. Since $x^d\geq1$ for small $x$, we obtain $R=O(s^d)$.
If $d>0$, the same inequality shows that $R(x)$ is Cauchy as $x\downarrow0$:
for $x,y<\eta$, its oscillation is at most $A\eta^d$. Let its limit be
$c$. Sending $x\downarrow0$ in \eqref{eq:profile-comparison} gives
$|R(y)-c|\leq Ay^d$. The constant $c$ belongs to $\mathcal P$.
Adding back the subtracted primitive proves the result in both cases.
\end{proof}
The exceptional order $-1$ is excluded only from this estimate for an
arbitrary remainder. The logarithmic monomial itself has already been
integrated explicitly. In the argument below the exceptional remainder
order is handled by weakening it to order $-2$ before integrating.

\subsection{Radial Laplace integrals and their Euler equation}
For $m\geq0$ and $j,k\in\mathbb N$, set
\[
 \omega(r)=\sqrt{m^2+r^2},\qquad
 M_{j,k}(s)=\int_0^\infty r^{j+2}\omega(r)^{k-1}
                                      e^{-s\omega(r)}\,dr,\quad s>0.
\]
The value of the integrand at $r=0$ is immaterial. For $r>0$,
$\omega(r)\geq r>0$. The integrands for $k=0$ and $k=1$ are respectively
bounded by $r^{j+1}e^{-sr}$ and $r^{j+2}e^{-sr}$. These are integrable:
scaling and repeated integration by parts give
\[
 \int_0^\infty r^q e^{-sr}\,dr=q!s^{-q-1}\quad(q\in\mathbb N).
\]
The boundary terms vanish at infinity by exponential decay and at zero
when $q>0$; the initial case is $\int_0^\infty e^{-sr}dr=s^{-1}$.
The identity $\omega^2=m^2+r^2$ now proves integrability for every $k$,
by induction in steps of two, and gives
\begin{equation}\label{eq:profile-moment-recur}
 M_{j,k+2}=m^2M_{j,k}+M_{j+2,k}.
\end{equation}
For a fixed $s>0$, differentiation of the integrand in a neighborhood
where the parameter is at least $s/2$ is dominated in absolute value by
the integrand for $M_{j,k+1}(s/2)$. Dominated differentiation gives
\begin{equation}\label{eq:profile-moment-deriv}
 M_{j,k}'(s)=-M_{j,k+1}(s).
\end{equation}

There is a second relation which improves the approximation order. For
$r>0$, differentiation in the radial variable gives
\[
 \frac{d}{dr}\bigl(r^{j+3}e^{-s\omega(r)}\bigr)
 =(j+3)r^{j+2}e^{-s\omega(r)}
       -s\frac{r^{j+4}}{\omega(r)}e^{-s\omega(r)}.
\]
The two terms are integrable by the bounds just proved. The primitive
vanishes at zero and infinity, including when $m=0$. Integration therefore
yields the Euler identity
\begin{equation}\label{eq:profile-euler}
 sM_{j+2,0}(s)=(j+3)M_{j,1}(s).
\end{equation}
Combining \eqref{eq:profile-moment-recur}--\eqref{eq:profile-euler} gives
\begin{equation}\label{eq:profile-weighted-deriv}
 \frac{d}{ds}\bigl(s^{j+3}M_{j,1}(s)\bigr)
                      =-m^2s^{j+3}M_{j,0}(s).
\end{equation}
Indeed the derivative on the left is
$(j+3)s^{j+2}M_{j,1}-s^{j+3}(m^2M_{j,0}+M_{j+2,0})$,
and the Euler identity cancels the first and last terms.

\begin{proposition}\label{prop:profile-moments}
Every $M_{j,k}$ has power--logarithm approximations.
\end{proposition}
\begin{proof}
The initial bound is $0\leq M_{j,0}(s)\leq(j+1)!s^{-j-2}$, so the zero
polynomial gives an approximation of order $b_0=-j-2$. Suppose $M_{j,0}$
has an approximation of order $b\geq b_0$. The right-hand side of
\eqref{eq:profile-weighted-deriv} then has an approximation of order
$j+3+b\geq1$. Lemma~\ref{lem:profile-primitive}, followed by division
by $s^{j+3}$, gives an approximation of order $b+1$ to $M_{j,1}$.
Since $M_{j,0}'=-M_{j,1}$, another application of the same lemma gives
an approximation to $M_{j,0}$ of order $b+2$, hence of order $b+1$,
provided $b+1\neq-1$. In the remaining case $b=-2$, weaken the
approximation to $M_{j,1}$ from order $-1$ to order $-2$. Integration
then gives precisely the required order $-1$ for $M_{j,0}$.

Induction from $b_0$ proves all larger orders, and weakening covers every
smaller integer order. The first half of the argument now gives every
order for $M_{j,1}$ as well: choose $b\geq b_0$ as large as needed.
Finally \eqref{eq:profile-moment-recur} and closure under addition and
scalar multiplication prove the assertion for all $k$, by induction in
steps of two, simultaneously for every $j$.
\end{proof}

\subsection{From modes to the logarithmic obstruction}
A mode consists of a mass $m\geq0$, integers $j,k\geq0$, and a continuous
real function $P$ on $\mathbb R^3$ homogeneous of degree $j$ under positive
scalings. Its Euclidean profile is
\begin{equation}\label{eq:profile-mode}
 K_{m,P,k}(s)=\int_{\mathbb R^3}
       \frac{\omega_m(p)^kP(p)}{2\omega_m(p)}e^{-s\omega_m(p)}\,dp,
 \qquad \omega_m(p)=\sqrt{m^2+|p|^2}.
\end{equation}
Let $d\sigma$ denote ordinary surface area measure on the unit sphere.
The maximum of $|P|$ on that sphere is finite. In polar coordinates,
$|P(r\theta)|\leq Cr^j$, so absolute integrability follows from that
of $M_{j,k}$. Fubini's theorem and homogeneity therefore give the exact
normalization
\begin{equation}\label{eq:profile-angular}
 K_{m,P,k}(s)=\frac12\left(\int_{S^2}P(\theta)\,d\sigma(\theta)\right)
                         M_{j,k}(s).
\end{equation}
A track is a finite product of mode profiles; the empty track is $1$.
By Proposition~\ref{prop:profile-moments} and
Lemma~\ref{lem:profile-algebra}, every track has power--logarithm
approximations. The same is true of the real part of a finite complex
linear combination of tracks, because all the individual tracks are real.

\begin{lemma}[No inverse logarithmic leading term]\label{lem:profile-log}
If a real function $f$ has an approximation of order zero and
\[
 \lim_{s\downarrow0}s^8\bigl(\log s^{-1}\bigr)^2f(s)=A\in\mathbb R,
\]
then $A=0$.
\end{lemma}
\begin{proof}
Write $f=P+O(1)$ with $P\in\mathcal P$. Since
$s^8(\log s^{-1})^2\to0$, the same weighted limit holds for $P$.
Restrict to $s_n=2^{-n}$. A monomial $c_{a,k}s^a(\log s)^k$ contributes
\[
 c_{a,k}\bigl(\log(1/2)\bigr)^{k+2}
             \bigl(2^{-(8+a)}\bigr)^n n^{k+2}.
\]
Combine equal pairs $(a,k)$ and discard zero coefficients. The remaining
finite sum has distinct pairs $(\lambda,q)=(2^{-(8+a)},k+2)$,
with $\lambda>0$ and $q\geq2$. Order them lexicographically and choose
the largest pair $(\lambda_*,q_*)$. If the sum is empty, its limit is
zero. If $\lambda_*<1$, every term tends to zero: writing
$\lambda=e^{-c}$ with $c>0$ reduces this to $n^qe^{-cn}\to0$.
For completeness, $e^{cn}\geq(cn)^{q+1}/(q+1)!$ proves that limit.

If $\lambda_*\geq1$, then $D_n=\lambda_*^n n^{q_*}\to\infty$.
Each smaller term divided by $D_n$ tends to zero. With a smaller base,
exponential decay dominates any remaining power of $n$ by the same
estimate; with equal base, the ratio is a negative power of $n$.
Thus the finite sum divided by $D_n$ tends to its nonzero leading
coefficient. But the sum has the finite limit $A$, and division by
$D_n\to\infty$ makes its limit zero, a contradiction. Only the first
case is possible, and it has $A=0$.
\end{proof}

\begin{theorem}\label{thm:profile-exclusion}
Let $T:(0,\infty)\to\mathbb C$ satisfy
\[
 \lim_{s\downarrow0}s^8\bigl(\log s^{-1}\bigr)^2T(s)=A,
 \qquad A>0.
\]
Then $T$ cannot agree on any interval $(0,\delta)$ with a finite complex
linear combination of tracks built from \eqref{eq:profile-mode}.
\end{theorem}
\begin{proof}
If such agreement held, the real part of that combination would have an
order-zero approximation. Taking real parts in the displayed limit and
using the agreement would contradict Lemma~\ref{lem:profile-log}.
\end{proof}
This proves the exclusion once the positive trace asymptotic has been
established for the constructed theory. It imposes no restriction on an
infinite sum or a continuous distribution of masses, and no such exclusion
is asserted here. The construction of $T$ and its positive coefficient $A$
is a separate quantitative step.
